\subsection{线性化}\label{sec:8i}
\eref{8.9}的解集 $\M_S(\{y_\zeta\})$ 可以置于标准的函数空间框架中。固定 $2<p<\infty$。设 $\mathcal B_S$ 是局部属于 $W^{1,p}$、并在条带状端上以 $W^{1,p}$ 意义收敛于某组 $\{y_\zeta\}$ 的映射 $u:S\to M$ 所成的 Banach 流形。这个流形在点 $u$ 处的切空间正是 $W^{1,p}(S,E,F)$，其中 $E=u^*TM$，而对每个分支 $C\subset\partial S$，有 $F|_C=u^*TL_C$。在 $E$ 上赋予殆复结构 $J_E(z)=J(z,u(z))$，这里 $J$ 是扰动数据的一部分。存在自然的 Banach 向量丛 $\mathcal E_S\to\mathcal B_S$，其纤维是空间 $L^p(S,\Omega_S^{0,1}\otimes u^*TM)$。可以把 $\overline\partial_J-\nu$ 看成这个丛的光滑截面，显然，其零点集就是各模空间 $\M_S(\{y_\zeta\})$ 的并；这也给这些空间赋予了适当的拓扑。模空间中一点处的线性化记为
\begin{equation}\label{eq:8.17}
D_{S,u}:(T\mathcal B_S)_u\longrightarrow(\mathcal E_S)_u,
\end{equation}
它是非退化 Cauchy–Riemann 算子，因而是 Fredholm 算子。
\begin{definition*}
若 $D_{S,u}$ 是满射，就称 $u$ \emph{正则}。若对所有 $u\in\M_S(\{y_\zeta\})$ 都如此，则称模空间本身\emph{正则}。这意味着它是一个适当维数的光滑流形，其维数等于 $D_{S,u}$ 的指标。若所有模空间都正则，
\SourcePage{116}
则称扰动数据 $(K,J)$ 本身\emph{正则}；对 Floer 数据也采用同样的术语。
\end{definition*}
为了将 $D_{S,u}$ 识别为前面所考虑的某个具体算子 $\overline\partial_\nabla$，在 $E$ 上赋予从 $\varpi$ 拉回的辛结构。联络 $\nabla$ 并不唯一，因为算子的定义只涉及它的 $(0,1)$ 分量；一种可能的选择方式如下。设 $\tilde u:S\to S\times M$ 是 $u$ 的图。固定 $Z\in C^\infty(S\times M,T^*S\otimes TM)$，满足 $Z(z,u(z))=du(z)$，且处处有 $(Z-Y)^{0,1}=0$。每个这样的 $Z$ 都对应 $E$ 上的一个联络 $\nabla$。具体地，取 $\widetilde X\in C^\infty(S\times M,TM)$，将其沿 $\tilde u$ 的拉回记为 $X\in C^\infty(S,E)$，令
\begin{equation}\label{eq:8.18}
\nabla X=\tilde u^*[Z,\widetilde X].
\end{equation}
为看出这个定义不依赖于 $\widetilde X$ 的选择，而且确实定义了联络，在 $TM\to S\times M$ 上任取无挠联络 $\widetilde\nabla^0$，将其沿 $\tilde u$ 的拉回记为 $\nabla^0$。借助它，可将\eref{8.18}改写为
\[
\nabla_\xi X=\nabla^0_\xi X-\tilde u^*(\widetilde\nabla^0_X\langle Z,\xi\rangle)
\]
对所有 $\xi\in C^\infty(S,TS)$ 成立；这里 $\langle Z,\xi\rangle\in C^\infty(S\times M,TM)$ 来自自然配对 $TS\otimes(T^*S\otimes TM)\to TM$。这表明 $\nabla$ 等于 $\nabla^0$ 加上一个零阶项，该项依赖于 $Z$ 沿 $u$ 的图的导数。接下来，用 $S$ 上的局部全纯坐标 $z=s+it$ 把计算写得更明确，得到 $\nabla$ 的 $(0,1)$ 部分为
\[
\begin{aligned}
(\nabla X)^{0,1}
={}&\tfrac12\bigl(\nabla^0_{\partial_s}X+J\nabla^0_{\partial_t}X
-\widetilde\nabla^0_X\langle Z,\partial_s\rangle-J\widetilde\nabla^0_X\langle Z,\partial_t\rangle\bigr)\\
={}&\tfrac12\bigl(\nabla^0_{\partial_s}X+J\nabla^0_{\partial_t}X
+\tilde u^*(\widetilde\nabla^0_XJ)\,du(\partial_t)\bigr)\\
&-\tfrac12\tilde u^*\widetilde\nabla^0_X\bigl(Y(\partial_s)+JY(\partial_t)\bigr).
\end{aligned}
\]
这确实是伪全纯映射理论中线性化算子的标准形式，最后一项是为计入\eref{8.9}中的非齐次项而加入的。因此 $D_{S,u}=\overline\partial_\nabla$。为了满足容许性条件，选择 $Z$，使得当 $s+it$ 是某个条带状端 $\epsilon_\zeta:Z^\pm\to S$ 上的标准坐标时，$Z(\partial_s)$ 在 $s\to\pm\infty$ 时以指数速率趋于零，收敛在 $C^2$ 拓扑中理解。此时，$(E,\varpi_E,J_E,F)$ 及 $\nabla$ 的渐近结构由以下极限数据描述：向量丛 $E_\zeta=y_\zeta^*TM$，带有从 $\varpi$ 继承的辛结构，以及殆复结构 $J_{E_\zeta}(t)=J_\zeta(t,y_\zeta(t))$；Lagrangian 子空间 $\Lambda_{\zeta,k}=T(L_{\zeta,k})_{y_\zeta(k)}$；以及沿积分曲线 $y_\zeta$ 线性化 $H_\zeta$ 的“流”所得到的辛联络 $\nabla_\zeta$。

\begin{note}\label{st:8.1}
在某些特殊情形，线性化算子有更直接的描述。假设 $J$ 不依赖于 $z\in S$ 且可积，并且 $Y=0$。那么 $TM$ 及其拉回 $E$ 自然是全纯丛，
\SourcePage{117}
$D_{S,u}$ 就是对应的典范 Dolbeault 算子。另一种情形是 $u\equiv x$ 为常值，而 $K$ 及其一阶导数在 $x$ 处均为零。此时二阶导数 $D^2K$ 可视为 $\Omega^1(S,\operatorname{sym}^2(T M_x)^\vee)$ 中的元素，因而给出平凡丛 $E=S\times T M_x$ 上的辛联络。于是 $D_{S,u}$ 是这个联络相对于复结构 $J_E(z)=J(z,x)$ 的 $(0,1)$ 部分。
\end{note}

\subsection{TQFT 的定义}\label{sec:8j}
对每一对正合 Lagrangian 子流形 $(L_0,L_1)$，选择正则 Floer 数据 $(H_{L_0,L_1},J_{L_0,L_1})$。
\begin{definition*}
定义 \emph{Floer 上链群} $CF^{\mathrm{pr}}(L_0,L_1)$ 为 $\K$-向量空间，对每个 $y\in\mathcal C(L_0,L_1)$ 赋予一个基向量。这个群上的边界算子为
\begin{equation}\label{eq:8.19}
\partial(y_1)=\sum_{y_0}\#\M_Z^*(y_0,y_1)\,y_0,
\end{equation}
其中 $\#$ 表示对孤立点的个数模 $2$ 计数。$HF^{\mathrm{pr}}(L_0,L_1)$ 是 $(CF^{\mathrm{pr}}(L_0,L_1),\partial)$ 的上同调。\znote{109}{这里有两处命名要留意。其一，$CF^{\mathrm{pr}}$ 虽称「上链群」，但 (8.19) 的 $\partial$ 把 $y_1$ 送到 $y_0$，按几何次数是降次的。其二，在 $\operatorname{char}(\K)=2$ 的初步理论里没有分次可谈，同调与上同调只差一次整体平移，命名不影响内容；到了带分次的精细理论（§9），次数由 Maslov 型指标给出，$\partial$ 降低该次数一次，此时必须固定约定，否则符号与次数会一起错位。[来源：原书 §8b、§8j 本段与 §9 关于分次的讨论。]}
\end{definition*}
接下来，对每个带有 Lagrangian 标记 $\{L_C\}$ 及条带状端的边界带标点 Riemann 曲面 $S$，选择与上述 Floer 数据相容的正则扰动数据 $(K_S,J_S)$。定义映射
\begin{equation}\label{eq:8.20}
\begin{aligned}
C\Phi_S:&\ \bigotimes_{\zeta^+\in\Sigma^+}CF^{\mathrm{pr}}(L_{\zeta^+,0},L_{\zeta^+,1})
\longrightarrow\bigotimes_{\zeta^-\in\Sigma^-}CF^{\mathrm{pr}}(L_{\zeta^-,0},L_{\zeta^-,1}),\\
C\Phi_S(\otimes_{\zeta^+}y_{\zeta^+})
={}&\sum_{\{y_{\zeta^-}\}}\#\M_S(\{y_{\zeta^-},y_{\zeta^+}\})\,(\otimes_{\zeta^-}y_{\zeta^-}),
\end{aligned}
\end{equation}
并以 $\Phi_S$ 表示它在上同调上诱导的映射。对这一领域不熟悉的读者应当记住：尽管这里的叙述相当简略，构造的每一步都是一项大工程。仅仅为了定义 $\partial$ 和 $C\Phi_S$，就需要：(i) 横截性结果，说明正则性是 Floer 数据与扰动数据的一般位置条件；(ii) 紧致性结果，保证在正则情形，相关模空间只有有限个孤立点。此后，适当的粘合定理说明：(iii) $\partial^2=0$；(iv) $C\Phi_S$ 是上链映射。进一步还需证明：(v) 所得理论在典范同构意义下不依赖于任何选择；(vi) 它满足复合公理\eref{8.2}。前四个问题在文献中已有充分论述，而且在构造 Fukaya 范畴时还会再次出现，见\secref{9k}、\secref{9l}，因此这里不再讨论。关于最后两个实际上相互交织的步骤，作一些说明是适当的。

首先着手 (v)。先证明，只要端上的 Floer 数据固定，$\Phi_S$ 就不依赖于 $S$ 的共形结构、条带状端的选择及扰动数据。为此，取由 $[0,1]$ 参数化的一族共形结构和扰动数据，在给定的两组之间插值，
\SourcePage{118}
并用所得全部模空间的并来定义一个上链同伦。这可以看作参数化模空间的首次出现，虽然是相当平凡的一例；后面会广泛讨论这种空间。接着转向 (vi)。取配备了条带状端、Lagrangian 边界条件及扰动数据的曲面 $S_1,S_2$，像\eref{8.5}那样沿无穷远点 $\zeta_1^+,\zeta_2^-$ 将它们粘合，形成 $S=S_1\#_lS_2$。假设这两个端的 Floer 数据相同，从而在 $S$ 上自然得到扰动数据。再假设满足适当的正则性条件，且 $l\gg0$，就可以证明复合律在\emph{上链层次}成立。当然，这个等式只对 $S$ 上这组非常特殊的共形结构及扰动数据成立。然而，已经证明，上同调层次的映射 $\Phi_S$ 不依赖于这些选择。对于 Fukaya 范畴这种上链层次的结构，这类论证将不再适用，这意味着我们必须对构造中所用的扰动数据施加精细得多的控制。

有了这些，(v) 中关于 Floer 数据选择无关性的剩余部分就完全是形式上的。考虑无限条带 $Z$，在边界分支 $C_k=\R\times\{k\}$（$k=0,1$）上分别赋予 Lagrangian 标记 $L_k$。在两个端 $\zeta_0=-\infty$、$\zeta_1=+\infty$，为对子流形 $(L_{\zeta_j,0},L_{\zeta_j,1})=(L_0,L_1)$ 配备正则 Floer 数据
\[
(H_{\zeta_j},J_{\zeta_j})=(H^j_{L_0,L_1},J^j_{L_0,L_1}),
\]
并选择相容的正则扰动数据。这给出所得 Floer 上同调群之间的映射
\[
\Phi_Z:HF^{\mathrm{pr}}(L_0,L_1)^1\longrightarrow HF^{\mathrm{pr}}(L_0,L_1)^0,
\]
通常称为\emph{延拓映射}\bcite{118}。上述同伦论证说明它是典范的；应用复合律则说明它是同构。

\begin{note}\label{st:8.2}
在许多情形，可以采用非齐次项为零的 Floer 数据与扰动数据。具体说，只要 $L_0\pitchfork L_1$，就能找到 $H=0$ 的正则 Floer 数据。类似地，考虑带有条带状端及 Lagrangian 边界条件 $\{L_C\}$ 的 Riemann 曲面 $S$。假设每个无穷远点所对应的一对 Lagrangian 子流形都横截相交，而且 $M$ 中没有同时位于所有 $L_C$ 上的点；第二个条件排除了\eref{8.9}的常值解。此时就能找到 $K=0$ 的正则扰动数据。
\end{note}

\subsection{移动边界条件}\label{sec:8k}
\begin{definition*}
设 $S$ 是带有条带状端的 Riemann 曲面。$S$ 的一组\emph{移动边界条件}，是随 $z\in\partial S$ 光滑变化的正合 Lagrangian 子流形族 $L_z\subset M$，并且当 $z\in\partial S\cap\epsilon_\zeta(Z^\pm)$ 时局部不变。因此，每个无穷远点仍有与之对应的一对子流形 $(L_{\zeta,0},L_{\zeta,1})$。
\end{definition*}
在扰动数据
\SourcePage{119}
$(K,J)$ 的定义中，应将\eref{8.6}替换为以下稍复杂的条件：
\begin{equation}\label{eq:8.21}
d(K(\xi)|_{L_z})=\varpi(\cdot,\partial_\xi L_z)\in\Omega^1(L_z)
\qquad\text{对任意 }\xi\in T(\partial S)_z.
\end{equation}
这里 $\partial_\xi L_z$ 是 $C^\infty(TM|_{L_z})$ 中任意一个表示该 Lagrangian 子流形无穷小变形的截面，因此\eref{8.21}右端是这种变形所对应的标准正合一形式。于是可以考虑与前面相同的方程\eref{8.9}。$K$ 的条件使得\eref{8.11}中的边界项可以通过积分处理掉，其余理论无需重大调整便可照常进行，得到推广的映射\eref{8.1}，满足 TQFT 公理的显然变体。\znote{111}{(8.21) 说的是：$L_z$ 随 $z$ 变动时，$K(\xi)$ 在 $L_z$ 上的限制不再是零（那是 (8.6) 的情形），而是一个恰当一形式，其外微分恰为 $\varpi(\cdot,\partial_\xi L_z)$，即该无穷小变形对应的标准正合一形式。这样 (8.11) 中沿边界的积分项可以写成恰当形式而被抵消，能量恒等式与 TQFT 公理才继续成立。紧接着的 (8.22) 又说明：只要各边界分支都是区间，移动边界条件其实可化为固定边界条件，故这一推广增加的是方便，而非实质内容。[来源：原书 (8.21)--(8.22) 及本段论证。]}

实际上，这一推广的幅度比初看起来小得多。为看清这一点，假设 $\partial S$ 的所有连通分支都是开区间。于是，存在由 $z\in S$ 参数化的自同构族 $\phi_z\in\operatorname{Aut}(M,\partial M)$，使得 $\phi_z(L_z)$ 局部不变，并且在条带状端上有 $(\partial/\partial s)\phi_{\epsilon_\zeta(s,t)}=0$。坐标变换
\begin{equation}\label{eq:8.22}
u(z)\longmapsto\phi_z(u(z))
\end{equation}
保持方程\eref{8.9}的类型，同时将移动边界条件化为固定的边界条件。尽管如此，用移动边界条件来思考往往很方便，以下应用说明了这一点。设 $\{L_s\}_{0\leq s\leq1}$ 是正合 Lagrangian 同痕。取 $S=H$，配备移动边界条件 $L_{\psi(s)}$，其中 $\psi\in C^\infty(\R,[0,1])$ 满足 $s\ll0$ 时 $\psi(s)=0$，$s\gg0$ 时 $\psi(s)=1$。由此得到一个特征元素
\begin{equation}\label{eq:8.23}
\Phi_H\in HF^{\mathrm{pr}}(L_0,L_1),
\end{equation}
它实际上就是\secref{8c}所提到的元素。那里陈述的全部性质，都由标准的同伦和粘合论证得到。例如，将\eref{8.23}与三角形乘积复合，对任意 $L$ 得到同构 $HF^{\mathrm{pr}}(L,L_0)\to HF^{\mathrm{pr}}(L,L_1)$ 及 $HF^{\mathrm{pr}}(L_1,L)\to HF^{\mathrm{pr}}(L_0,L)$。事实上，细察粘合论证并应用变换\eref{8.22}，就可将这些同构归结为先前考虑的延拓映射。

\subsection{边界点处的求值}\label{sec:8l}
假设给定一个带有 Lagrangian 标记的边界带标点 Riemann 曲面，并给定额外的有限标点集 $\Pi\subset\partial S$，分为负标点 $\Pi^-$ 和正标点 $\Pi^+$。此时可以引入 TQFT 映射\eref{8.1}的推广：
\begin{equation}\label{eq:8.24}
\begin{aligned}
\Phi_{S,\Pi}:&\ \bigotimes_{\zeta^+\in\Sigma^+}HF^{\mathrm{pr}}(L_{\zeta^+,0},L_{\zeta^+,1})
\otimes\bigotimes_{z^+\in\Pi^+}H^*(L_{C_{z^+}};\K)\\
&\longrightarrow\bigotimes_{\zeta^-\in\Sigma^-}HF^{\mathrm{pr}}(L_{\zeta^-,0},L_{\zeta^-,1})
\otimes\bigotimes_{z^-\in\Pi^-}H^*(L_{C_{z^-}};\K).
\end{aligned}
\end{equation}
\SourcePage{120}
这里 $C_z$ 是 $z$ 所属的边界分支。除了通常将无穷远点粘合外，现在还可以作边界连通和：在相应 Lagrangian 边界条件相同的点 $z_1^+\in\Pi_1^+$、$z_2^-\in\Pi_2^-$ 处，把两个曲面 $S_1,S_2$ 接合起来。对于对应的算子\eref{8.24}，这相当于在相关因子 $H^*(L_{C_z};\K)$ 上施行复合。

最简单的情形是 $S$ 紧致。假设正则性成立，则\eref{8.9}的解所成的模空间 $\M_S$ 是光滑闭流形。求值映射 $\operatorname{ev}_z(u)=u(z)$（$z\in\Pi$）的乘积，在 $\prod_zL_{C_z}$ 中给出一个同调类。利用 Künneth 公式以及针对 $\Pi^-$ 因子的 Poincaré 对偶，可将它解释为多重线性映射\eref{8.24}。然而，如果令非齐次项为零，直接的能量论证就表明，$\M_S$ 只包含常值映射，因而所得不变量并没有多少新意。\znote{112}{这一段把 $H^*(L)$ 也纳入 TQFT：额外的边界标点 $z$ 上可以插入 $L_{C_z}$ 的上同调类，几何上就是把解在 $z$ 处的取值与 $\prod_zL_{C_z}$ 中的闭链求交。$S$ 紧致时，求值映射的乘积给出 $\prod_zL_{C_z}$ 中的同调类，用 Künneth 公式与 Poincaré 对偶即可读出 (8.24)。但若取零非齐次项，能量论证迫使解为常值映射，于是只剩 (8.25) 列出的四种初等映射；真正有用的是带一个出向端与一个边界标点的情形：取 $S=\overline H$ 就得到 PSS 同构 (8.4)。[来源：原书 §8l 本段与 (8.24)--(8.25)，并见原书文献 [120]。]}例如，取 $S=D$，Lagrangian 标记为 $L$，得到
\begin{equation}\label{eq:8.25}
\begin{cases}
\text{单位元 }1\in H^*(L;\K),&|\Pi^-|=1,\ |\Pi^+|=0,\\
\text{积分映射 }H^*(L;\K)\to\K,&|\Pi^-|=0,\ |\Pi^+|=1,\\
\text{恒等映射 }H^*(L;\K)\to H^*(L;\K),&|\Pi^-|=|\Pi^+|=1,\\
\text{杯积},&|\Pi^-|=1,\ |\Pi^+|=2.
\end{cases}
\end{equation}
为简要说明更一般的情形，假设 $\Sigma=\{\zeta\}$ 只含一个出向端，$\Pi=\{z\}$ 只含一个负标点。其思路是：取 Floer 上闭链 $\sum_km_ky_k\in CF^{\mathrm{pr}}(L_{\zeta,0},L_{\zeta,1})$，则对应的求值映射 $\operatorname{ev}_z:\M(y_k)\to L_{C_z}$ 的形式线性组合，在通常的拓扑意义下是 $L_{C_z}$ 上的闭链。其同调类，更准确地说是它的 Poincaré 对偶，就是 Floer 上同调类 $[\sum_km_ky_k]$ 在\eref{8.24}下的像。

一种等价而在技术上可能更简洁的做法，可以追溯到\bcite{120}，即采用 $H^*(L_{C_z};\K)$ 的 Morse 同调模型。这意味着，\eref{8.24}底层上链映射的系数，通过计数交集 $\M(y_k)\cap\operatorname{ev}_z^{-1}(W^u(x))$ 中的点来确定；这里 $W^u(x)$ 是 $-\nabla h_z$ 的不稳定流形，$h_z$ 是在 $L_{C_z}$ 上选取的 Morse 函数。例如，取 $S=\overline H$，Lagrangian 标记为 $L$，便可定义 PSS 映射\eref{8.4}。\secref{8c}中陈述的所有性质，现在都由\eref{8.25}及复合律推出。例如，逆映射恰由带有一个正标点的 $S=H$ 给出。从某种意义上说，这实际上是唯一需要的应用：可以利用粘合及 PSS 同构，把标点 $z$ 替换为满足 $L_{\zeta,0}=L_{\zeta,1}$ 的无穷远点 $\zeta$，从而将\eref{8.24}化为更基本的版本\eref{8.1}。
